\documentclass{article}
\usepackage[utf8]{inputenc}
\usepackage{hyperref}
\usepackage{listings}
\usepackage{xcolor}
\usepackage{enumitem}

\lstset{
    basicstyle=\ttfamily\small,
    breaklines=true,
    frame=single,
    backgroundcolor=\color{gray!5},
}

\title{Replace string in multiple files using sed}
\author{Marcos de Carvalho}
\date{2026-04-21}

\begin{document}

\section*{Replace string in multiple files using sed}
\textbf{Author:} Marcos de Carvalho \hfill \textbf{Date:} 2026-04-21

\noindent Replaces all occurrences of a string with a replacement string in multiple files or file patterns. Pure oneliner that works directly pasted with parameters, or in an alias.

\section*{Language}
bash

\section*{Category}
text-processing

\section*{Command}
\begin{lstlisting}
bash -c 'old=$1; new=$2; shift 2; find "$@" -type f -exec sed -i "s#$old#$new#g" {} +' _
\end{lstlisting}

\section*{Explanation}
This oneliner uses a subshell (`bash -c`) to parse positional parameters. The first argument is the string to replace (`\$1`), the second is the replacement (`\$2`), and `shift 2` leaves the remaining arguments as the files or patterns to search in (`\$@`). The `find` command locates all matching files and processes them with `sed` using the `-exec \{\} +` syntax, which efficiently handles multiple files in batches. The `\_` at the end acts as `\$0` (the script name) so `\$1` maps to the first actual argument. This approach handles an arbitrary amount of files and supports wildcard file patterns without hardcoding values into the command itself.

\section*{Tags}
sed, text-processing, string-replacement, file-editing, search-replace, batch-replace

\section*{Dependencies}
sed, coreutils


\section*{Arguments}
\begin{enumerate}
  \item \textbf{SEARCH} (Required): The string to search for (first parameter passed after the command).
  \item \textbf{REPLACE} (Required): The replacement string (second parameter passed after the command).
  \item \textbf{FILES} (Required): One or more file paths or patterns (e.g., file1.txt, *.txt, /path/*.js).
\end{enumerate}

\section*{Examples}
\begin{enumerate}
  \item \texttt{bash -c 'old=\$1; new=\$2; shift 2; find "\$@" -type f -exec sed -i "s\#\$old\#\$new\#g" \{\} +' \_ 'old-domain.com' 'new-domain.com' config.json settings.json} - Paste directly: replace 'old-domain.com' with 'new-domain.com' in config.json and settings.json
  \item \texttt{bash -c 'old=\$1; new=\$2; shift 2; find "\$@" -type f -exec sed -i "s\#\$old\#\$new\#g" \{\} +' \_ '/var/old/path' '/var/new/path' nginx.conf} - Replace a path containing slashes in a configuration file
  \item \texttt{bash -c 'old=\$1; new=\$2; shift 2; find "\$@" -type f -exec sed -i "s\#\$old\#\$new\#g" \{\} +' \_ 'old-text' 'new-text' *.txt} - Replace in all .txt files in current directory using a wildcard
  \item \texttt{alias repl="bash -c 'old=\textbackslash\{\}\$1; new=\textbackslash\{\}\$2; shift 2; find \textbackslash\{\}"\textbackslash\{\}\$@\textbackslash\{\}" -type f -exec sed -i \textbackslash\{\}"s\#\textbackslash\{\}\$old\#\textbackslash\{\}\$new\#g\textbackslash\{\}" \{\} +' \_"} - Create an alias named 'repl' to easily perform string replacements
  \item \texttt{repl 'deprecated\_method' 'new\_method' src/**/*.py} - Use the alias to replace in all Python files in src directory and subdirectories
  \item \texttt{repl 'old-api' 'new-api' config/*.json docs/*.md} - Use the alias to replace in multiple file types across different directories
\end{enumerate}

\section*{Output}
\begin{lstlisting}
(No output on success - files are modified in place)
\end{lstlisting}

\section*{Notes}
\begin{itemize}
  \item Pure oneliner: uses bash -c to accept arbitrary parameters after the command string
  \item Usage: Paste the command and append your search string, replacement string, and file patterns
  \item Can be easily turned into a powerful alias in your .bashrc or .bash\_profile
  \item The -i flag modifies files in place; make backups first with cp file.txt file.txt.bak if concerned
  \item Using \# as delimiter prevents issues when search/replace strings contain forward slashes
  \item The g flag replaces all occurrences per line; remove it to replace only first occurrence
  \item The -exec \{\} + syntax bundles multiple files to sed for efficiency
  \item Files with spaces or special characters are handled correctly by find
  \item For case-insensitive replacement, modify the command to include the i flag: sed -i "s\#\$old\#\$new\#gi"
\end{itemize}

\section*{Warnings}
\begin{itemize}
  \item The -i flag modifies files directly without confirmation - always test on copies first
  \item If search string is empty, sed will replace every character with the replacement string
  \item The command will silently fail to modify files if they lack write permissions
  \item Very large files or many files may cause performance issues
  \item Replacement is literal - regex special characters will be treated as patterns (not escaped)
  \item Be careful with wildcard patterns - they may match more files than intended
  \item The find command only processes regular files, not directories
  \item Special characters in search/replace strings must be properly quoted or escaped
\end{itemize}

\section*{See Also}
\begin{itemize}
  \item find-large-files-recursive
  \item disk-space-usage-per-directory
\end{itemize}

\section*{Status}
reviewed

\section*{Safety}
caution

\section*{Shell}
bash

\section*{Platforms}
linux, gnu-linux, freebsd, openbsd, netbsd

\section*{Created At}
2026-04-21

\section*{Updated At}
2026-04-21

\end{document}
